\documentclass[11pt]{article}
\usepackage{mathblatt}
% Präambel des Setzers (werkzeuge/setzer-praeambel.tex, wortgleich); dazu \gteil (Muster T2U)
\makeatletter
\setlength{\mbpfbe}{0mm}% keine Punkte (pruefung.md § 5: nur beim Prüfstein)
\newlength{\gEin}\setlength{\gEin}{0pt}
\let\gSchrift\relax
\renewcommand{\pfkreuzzeile}[1]{\par\noindent\hangindent3.2em\hangafter1\hspace*{1.6em}$\square$\ #1\par}
\newcommand{\gkreuz}[1]{$\square$\ #1\hspace{2em}}
\newcommand{\gkreuzzeile}[2]{\par\noindent\hangindent2.2em\hangafter1\mbox{#1\,$\square$}\ #2\par}% Buchstabe, Option: eine je Zeile (neben Skizze)
% links, was man ansieht, rechts, was man tut; Spaltengrenze überall an derselben Stelle (gemeinsam.md „Satz“)
% v0.17: Fußeintrag „Nr: Wert“ – kurz in einer Box (nie getrennt), lang umbrechbar mit der Nummer am Anfang
% (vorher \mbox: lange Lösungen liefen über den Rand, AS40 138 pt); robust, weil \fusshilfe in die .aux schreibt
\NewDocumentCommand{\gfuss}{m m}{\sbox0{#1:~#2}\ifdim\wd0<0.45\textwidth\mbox{#1:~#2}\else\mbox{#1:}\nobreak\ #2\fi}
\newsavebox{\gbox}\newlength{\gLinks}
% v0.10 (Lehrer 13.10.): keine starre Halbierung – die linke Spalte ist so breit wie die Figur (+ 6 mm), die
% Antwort steht direkt daneben; ist die Figur breiter als 0,62 Zeile, steht die Antwort darunter
\newcommand{\gzwei}[2]{\par\smallskip\noindent\sbox{\gbox}{#1}%
  \setlength{\gLinks}{\dimexpr\wd\gbox+6mm\relax}%
  \ifdim\gLinks>0.62\linewidth
    \usebox{\gbox}\par\smallskip\noindent\begin{minipage}{\linewidth}#2\end{minipage}%
  \else
    \begin{minipage}[t]{\gLinks}\vspace{0pt}\usebox{\gbox}\end{minipage}%
    \begin{minipage}[t]{\dimexpr\linewidth-\gLinks\relax}\vspace{0pt}\raggedright #2\end{minipage}%
  \fi\par}
% v0.9 (12.10., pruefung.md § 5 „Marke mit B, Bedingungen rechts“): drittes Argument von pfaufg = Bedingungen
% der Prüfung (v0.10: \pfohnetr = TR durchgestrichen, „FOR“), klein grau im rechten Rand auf der Grundlinie der Nummer, Zeilen mit \\
\newcommand{\gbedrand}[2]{\ifblank{#2}{}{\rlap{\hspace*{\dimexpr#1+2mm\relax}%
  \parbox[t]{15mm}{\raggedright\scriptsize\leavevmode\color{mbgrau}\strut#2}}}}
% Teil eines Bündels (a, b …): Marke links im Rand, Bedingungen rechts im Rand, auf der Zeile des Teils
\newcommand{\gteilrand}[2]{\llap{\mbpfrandmarke{#1}\hspace*{\mbpfnr}}\gbedrand{\linewidth}{#2}}
% Teil eines Bündels: Buchstabe in eigener Spalte, Inhalt eingerückt (\gEin = Einzug, damit die Spaltengrenze von
% \gzwei an derselben Stelle bleibt)
\newenvironment{gbteil}[3]{\par\smallskip\noindent\gteilrand{#2}{#3}\makebox[\mbpfteil][l]{\textbf{#1}}%
  \begin{minipage}[t]{\dimexpr\linewidth-\mbpfteil\relax}\setlength{\gEin}{\mbpfteil}\raggedright\strut\ignorespaces}%
 {\end{minipage}\par\addvspace{6pt}}% v0.18 (PG80 Nr. 1): Abstand, damit die Linie von a) nicht an b) rutscht
% Grafik höchstens so breit wie die Spalte (Bündel nebeneinander, Koordinatensystem ohne Karo)
\newsavebox{\gpbox}
\newcommand{\gpassend}[1]{\sbox{\gpbox}{#1}\ifdim\wd\gpbox>\linewidth\typeout{GPASSEND: verkleinert \the\wd\gpbox\space auf \the\linewidth}\resizebox{\linewidth}{!}{\usebox{\gpbox}}\else\usebox{\gpbox}\fi}
% Skizzenfeld (v0.10, Lehrer 13.10.): Karo ab dem linken Rand über die Zeile, „Skizze“ klein grau oben links im Feld
% (wie \pfskizzenfeld der Vorlage); der rechte Rand bleibt den Bedingungen
\newcommand{\gskizze}[1]{\par\smallskip\noindent\begin{tikzpicture}
  \clip (0,0) rectangle ({\linewidth-1mm},{#1*10mm});
  \draw[black!16,line width=0.3pt,step=5mm] (0,0) grid ({\linewidth},{#1*10mm});
  \node[anchor=north west,inner sep=2pt,font=\scriptsize,text=mbgrau,fill=white] at (0,{#1*10mm}) {Skizze};
  \end{tikzpicture}\par}
\RenewDocumentEnvironment{pfaufg}{m m m}{%
  \begin{lrbox}{\mbaufgabebox}\begin{minipage}[t]{\linewidth}%
  \gSchrift\noindent\hspace*{\gEin}\mbpfrandmarke{#1}\gbedrand{\linewidth-\gEin}{#3}%
  \begin{minipage}[t]{\mbpfnr}\strut\textbf{#2}\end{minipage}%
  \begin{minipage}[t]{\dimexpr\linewidth-\gEin-\mbpfnr\relax}\raggedright\strut\ignorespaces}%
 {\end{minipage}%
  \end{minipage}\end{lrbox}%
  \setlength{\mbaufgabehoehe}{\dimexpr\ht\mbaufgabebox+\dp\mbaufgabebox\relax}%
  \ifdefempty{\gkopfnach}{\par\addvspace{5pt}\Needspace*{\mbaufgabehoehe}}{%
    % Stufenkopf, der mit dieser (ersten) Aufgabe zusammenbleibt: Platz für beide verlangen, dann den Kopf setzen
    % (12.10., Quadratische: ein Kopf stand allein unten, die Aufgabe mit Koordinatensystem auf der nächsten Seite)
    % v0.17 (AS40 „Terme deuten“ allein unten): hinter dem Kopf steht das Sprungziel (\special, ein Knoten, nach
    % dem die Seite brechen durfte) – darum \nobreak nach dem Kopf; Vorhalt 56 pt = Kopf mit Linie und Abständen
    \par\Needspace*{\dimexpr\mbaufgabehoehe+56pt\relax}\gkopfnach\par\nobreak\gdef\gkopfnach{}\addvspace{5pt}}%
  \noindent\usebox{\mbaufgabebox}\par\addvspace{5pt}}
% v0.18 (plan.md § 10, 15.10. l): Abschnittsbuchstabe A, B, C vor dem Namen im Kopf – als Vorsatz, damit der
% Name im TeX gleich bleibt (Bestellkarte, Sprungziele lesen ihn)
\def\gbstakt{}\newcommand{\gbstvor}[1]{\gdef\gbstakt{#1}}
\NewCommandCopy\gpfstufekopfalt\pfstufekopf
\RenewDocumentCommand{\pfstufekopf}{o m m}{\IfValueTF{#1}{\gpfstufekopfalt[#1]}{\gpfstufekopfalt}%
  {\ifx\gbstakt\empty\else\gbstakt\enspace\fi #2}{#3}\gdef\gbstakt{}}
% v0.18 (16.10. e): Begriff grau, ohne Nummer, ohne Rahmen über der ersten Aufgabe der Stufe
\newcommand{\gbegriff}[1]{\par\nobreak\noindent\hspace*{\mbpfnr}\parbox[t]{\dimexpr\linewidth-\mbpfnr\relax}{\raggedright\small\color{mbgrau}#1}\par\nobreak\addvspace{3pt}}
% v0.18 (16.10. g): Lösungsblatt – Kopf mit kräftigem Strich, Abschnitt klein grau als Zwischenzeile
\newcommand{\glkopf}[1]{\noindent{\large\bfseries Lösungen \textperiodcentered{} #1}\par\vspace{3pt}\hrule height 1.2pt\par\medskip}
\newcommand{\glzwischen}[1]{\par\addvspace{8pt}\Needspace*{4\baselineskip}\noindent{\footnotesize\color{mbgrau}#1}\par\nobreak\vspace{1pt}}
% v0.18 (15.10. o, w): Blattende – Kern ohne Rahmen, dann graue Zeile „Mehr: n Aufgaben“
\newcommand{\gkern}[1]{\par\addvspace{14pt}\Needspace*{6\baselineskip}\noindent{\bfseries Kern}\par\nobreak\smallskip
  {\small\raggedright #1\par}}
\newcommand{\gmehr}[1]{\par\addvspace{10pt}\noindent{\footnotesize\color{mbgrau}Mehr: #1 Aufgaben}\par}
% Stufenkopf vormerken: \gkopfvor{<Kopf>} – die nächste pfaufg setzt ihn, sobald beide zusammen Platz haben
\def\gkopfnach{}
\newcommand{\gkopfvor}[1]{\gdef\gkopfnach{#1}}
% Teilaufgabe mit Skizze links in gzwei: Buchstabe in eigener Spalte, darunter Text und Skizze
\newcommand{\gteil}[2]{\makebox[\mbpfteil][l]{\textbf{#1}}\parbox[t]{\dimexpr0.5\textwidth-\gEin-\mbpfnr-5mm-\mbpfteil\relax}{\raggedright #2}}
% Stufenkopf der Vorlage (mathblatt.sty 2026-10-12: Linie, Abstand, 13 pt, \pfjahre) ohne ihren Vorhalt von
% 14 Zeilen: der Setzer fordert vor dem Kopf 0,16 Seitenhöhe (Kopf + erste Aufgabe), mehr nicht
\setlength{\mbpfstufevorhalt}{0pt}
\makeatother
\newcommand{\karo}[1]{\par\smallskip\noindent\begin{tikzpicture}
  \clip (0,0) rectangle ({\linewidth-1mm},{#1*10mm});
  \draw[black!16,line width=0.3pt,step=5mm] (0,0) grid ({\linewidth},{#1*10mm});
  \end{tikzpicture}\par}
\newcommand{\kk}[1]{$\square$\,#1\hspace{0.9em}}
\newcommand{\linie}[1]{\,{\color{black!45}\rule[-1pt]{#1}{0.4pt}}\,}
\newcommand{\pkt}{\textperiodcentered{}}
% v0.14 (14.10.): Stufenkopf mit fünf Punkten (Option = Jahre für die Bestellkarte) oder „zuletzt 19“
\NewDocumentCommand{\gpunkte}{o m}{#2}
\newcommand{\gzuletzt}[1]{zuletzt #1}
% v0.7 (12.10.): Schülerfrage unter dem Stufenkopf; Ankreuzoption als Bild (Buchstabe und Kästchen darunter)
\newcommand{\gfrage}[1]{\par\nobreak\noindent\hspace*{\mbpfnr}{\small\itshape #1}\par\nobreak\addvspace{2pt}}
\newlength{\gbildbreite}\setlength{\gbildbreite}{0.32\linewidth}% Setzer setzt sie je Reihe (vier Bilder: 0,24)
\newcommand{\gbildoption}[3]{\begin{minipage}[t][#3][b]{\gbildbreite}\centering #2\par\smallskip\mbox{#1\,$\square$}\end{minipage}}
% Lineare (12.10.): Bild mit Buchstabe ohne Kästchen (Diagramm zuordnen); Aussage mit „wahr / falsch“ zum Ankreuzen
\newcommand{\gbildname}[3]{\begin{minipage}[t][#3][b]{\gbildbreite}\centering #2\par\smallskip\mbox{\textbf{#1}}\end{minipage}}
\newcommand{\gwahrfalsch}[2]{\par\noindent\hangindent2.2em\hangafter1\makebox[2.2em][l]{(#1)}#2\hfill\mbox{wahr $\square$\quad falsch $\square$}\par}
\makeatletter\setlength{\mbpfbe}{0mm}\makeatother
\begin{document}

\pfheftstil
\pflaufkopf[Abitur GK \pkt{} mehr]{Extrempunkte}{AA50}\pfseitevon
\pfheftkopf{Extrempunkte}{mehr}
\pfstufenzeile{Extrempunkte gesucht \pkt{} Hochpunkt mit e-Funktion gesucht \pkt{} Extremstelle nachweisen}
% 2018-be-gk-B1.1c-m1 (Schwere II, rang 11)
\begin{pfaufg}{nach\\GK~’18\\B1.1c}{1.}{}
Der Aufsprunghang einer Skisprungschanze wird beschrieben durch $g(x) = \tfrac{1}{1000} \cdot \left(\tfrac{1}{1600}x^4 - 8x^2 + 25\,600\right)$; 1 LE entspricht 1 m. Der Hang beginnt im Punkt $C(0 \mid 25{,}6)$; der Punkt $U$ liegt an der tiefsten Stelle. Es gilt $g'(x) = \tfrac{1}{1000} \cdot \left(\tfrac{1}{400}x^3 - 16x\right)$. Bestimme Lage und Art aller lokalen Extrempunkte des Graphen von $g$ und entscheide, welcher davon $U$ ist.
\karo{8}
\fusshilfe{\gfuss{1}{$-$80, 0, 80; Tiefpunkte ($-$80|0) und (80|0); U(80|0)}}
\end{pfaufg}
% 2018-be-gk-B1.2d-m1 (Schwere II, rang 12)
\begin{pfaufg}{GK~’18\\B1.2d}{2.}{}
Das Höhenprofil eines Wanderwegs wird für $0 \le x \le 8$ beschrieben durch $f(x) = (x + 1) \cdot e^{-\tfrac{x}{3}}$; 1 LE entspricht 1 km. Es gilt $f'(x) = \tfrac{2 - x}{3} \cdot e^{-\tfrac{x}{3}}$. Berechne die Koordinaten des höchsten Punkts des Profils. Die hinreichende Bedingung musst du nicht prüfen.
\karo{6}
\fusshilfe{\gfuss{2}{2; höchster Punkt (2|1,54): 2 km nach dem Start, 1,54 km hoch}}
\end{pfaufg}
% 2019-be-gk-B2.2b-m1 (Schwere II, rang 13)
\begin{pfaufg}{GK~’19\\B2.2b}{3.}{}
Der Spiegel eines Hormons im Blut wird in Prozent des Sollwerts beschrieben durch $h(t) = 6t \cdot e^{-0{,}05t} + 40$, $t \ge 0$; $t$ ist die Zeit in Tagen seit Beginn einer Behandlung. Es gilt $h'(t) = (6 - 0{,}3t) \cdot e^{-0{,}05t}$. Berechne, wann der Hormonspiegel am größten ist, und gib diesen Wert an. Die notwendige Bedingung genügt.
\karo{5}
\fusshilfe{\gfuss{3}{20; 84,1 \% nach 20 Tagen}}
\end{pfaufg}
% 2020-be-gk-B2.2c-m1 (Schwere II, rang 14)
\begin{pfaufg}{GK~’20\\B2.2c}{4.}{}
Gegeben ist die Funktion $f$ mit $f(x) = -0{,}02 \cdot (x - 6) \cdot (x + 2)^2 = -\tfrac{1}{50}x^3 + \tfrac{1}{25}x^2 + \tfrac{2}{5}x + \tfrac{12}{25}$, $x \in \mathbb{R}$. Bestimme Lage und Art der Extrempunkte des Graphen von $f$.
\karo{8}
\fusshilfe{\gfuss{4}{$\tfrac{10}{3}$; Tiefpunkt ($-$2|0); (3,33|1,52)}}
\end{pfaufg}
% 2021-be-gk-A1.3a-m1 (Schwere I, rang 15)
\begin{pfaufg}{GK~’21\\A1.3a}{5.}{}
Gegeben ist die Funktion $f$ mit $f(x) = x^3 - 6x^2 + 32$. Weise nach, dass $x_E = 4$ eine der beiden Extremstellen von $f$ ist.
\karo{6}
\fusshilfe{\gfuss{5}{0; Extremstelle (Tiefpunkt)}}
\end{pfaufg}
% 2021-be-gk-B2.1c-m1 (Schwere II, rang 16)
\begin{pfaufg}{GK~’21\\B2.1c}{6.}{}
Gegeben ist die Funktion $f$ mit $f(x) = \tfrac{1}{8}x^3 - x^2 + 2x$, $x \in \mathbb{R}$; ihr Graph ist $G_f$. Weise nach, dass $H\left(\tfrac{4}{3} \mid \tfrac{32}{27}\right)$ Hochpunkt von $G_f$ ist.
\karo{7}
\fusshilfe{\gfuss{6}{0; $-$1 < 0; H($\tfrac{4}{3}$|$\tfrac{32}{27}$) ist Hochpunkt}}
\end{pfaufg}
% 2021-be-gk-B2.2d-m1 (Schwere II, rang 17)
\begin{pfaufg}{GK~’21\\B2.2d}{7.}{}
Gegeben ist die in $\mathbb{R}$ definierte Funktion $f$ mit $f(x) = \left(-\tfrac{1}{5}x^2 + 3x\right) \cdot e^{-0{,}2x}$; ihr Graph ist $G$. Ohne Nachweis darfst du $f'(x) = \tfrac{1}{25} \cdot (x^2 - 25x + 75) \cdot e^{-0{,}2x}$ verwenden. Bestimme Koordinaten und Art der lokalen Extrempunkte von $G$ (auf zwei Nachkommastellen).
\karo{8}
\fusshilfe{\gfuss{7}{3,49 oder 21,51; (3,49|4,00); (21,51|$-$0,38)}}
\end{pfaufg}
% 2026-iqb-ga-A1.1a-m1 (Schwere I, rang 18)
\begin{pfaufg}{IQB~’26\\A1.1a}{8.}{}
Gegeben ist die in $\mathbb{R}$ definierte Funktion $f$ mit $f(x) = x^3 - 12x$; ihr Graph ist $G$. Weise nach, dass $G$ einen Hochpunkt mit der $x$-Koordinate $-2$ hat.
\karo{6}
\fusshilfe{\gfuss{8}{0; $-$2}}
\end{pfaufg}
% 2020-iqb-ga-A1.1a-m1 (Schwere I, rang 19)
\begin{pfaufg}{IQB~’20\\A1.1a}{9.}{}
Gegeben ist die in $\mathbb{R}$ definierte Funktion $f$ mit $f(x) = x^3 - 27x + 54$. Weise nach, dass $-3$ und $3$ die Extremstellen von $f$ sind.
\karo{5}
\fusshilfe{\gfuss{9}{3; 18 > 0 (Tiefpunkt)}}
\end{pfaufg}
\par\vfill\noindent{\scriptsize\color{mbgrau}Originale: Abi GK 2018 \pkt{} B1.1c, Abi GK 2018 \pkt{} B1.2d, Abi GK 2019 \pkt{} B2.2b, Abi GK 2020 \pkt{} B2.2c, Abi GK 2021 \pkt{} A1.3a, Abi GK 2021 \pkt{} B2.1c, Abi GK 2021 \pkt{} B2.2d, IQB GK 2026 \pkt{} A1.1a, IQB GK 2020 \pkt{} A1.1a\par}
\end{document}